\chapter{高斯表示与原语设计}

% CODEx_AUTHORED_CHAPTER: draft
\paragraph*{本章阅读锚点}
原始 3DGS 的基本原语是各向异性 3D Gaussian：它易投影、易排序、易 alpha 混合，也能通过密度控制适配复杂场景。但这个原语有明显边界：高斯核无限延伸且形状光滑，难以表达硬边、薄面、纹理局部变化、锐利高光、远近尺度统一和明确表面。因此本章的方法都在改“一个 primitive 到底是什么”：有的把高斯核换成 Student-t、半高斯、Beta、多项式或凸体；有的保留高斯几何但给它纹理或神经外观；有的把方向外观重新参数化；有的把 splatting 推到 2D 图像表示或 triangle/surfel 表面表示。

\section{替代核函数与新型基元}

替代核路线的锚点是原始 3DGS 的指数高斯核。它的优点是仿射投影有解析形式，缺点是对不连续边界、硬表面、厚尾覆盖和有限支持并不总合适。下面这些方法不是在高斯外面加一个 loss，而是直接替换或推广 splat 的数学形状。

% METHOD: 2406.02720
\paragraph{3D-HGS（2024）\cite{2406_02720}}
3D-HGS针对对称 Gaussian 在物体轮廓两侧同时扩散的缺陷。标准核在几何或颜色突变附近不可避免地把一侧贡献泄漏到另一侧，导致边缘模糊、过冲或需要大量小高斯补偿。该方法保留 Gaussian 的椭球投影、深度排序和可微 rasterizer 框架，但将每个基元的径向核换成有方向截断的 half-Gaussian。

半高斯在局部坐标中用一个可学习方向/分界面将核的一侧保留、另一侧抑制，使同一各向异性基元可以在平滑方向仍有连续衰减、在轮廓法向呈现单侧支撑。投影后，二维足迹与 alpha 计算要同时考虑原协方差和这个半空间条件；在像素位于有效一侧时使用核贡献，另一侧的 opacity 被压低或为零。这样可让高斯靠近边界但不跨边界混色，而不是先让多个对称 splat 重叠、再靠透明度拟合截断。

其余训练仍以渲染 RGB 与训练图的标准图像损失反传，位置、尺度、旋转、颜色与半核方向共同优化；alpha 合成路径没有被替换。半核是局部不连续先验，不自动解决遮挡、深度排序或未观测区域，且分界方向/截断选择若与真实表面不一致会引入孔洞或偏置。相比厚尾、负核等改变整体衰减形状的路线，3D-HGS的主要表达增益是单基元的非对称边界。

% METHOD: 2411.14974
\paragraph{3D Convex Splatting（2024）\cite{2411_14974}}
3D Convex Splatting用可微 smooth convex 替代高斯椭球。Gaussian 的等值面天然光滑且局部各向异性，表达平面、硬边和稠密块状体时常需叠加大量基元；一个 convex 则由一组可学习三维点定义，可在少量原语内覆盖更复杂的凸形状。

每个 primitive 的点集在相机下投影到二维，先构造投影点的凸包；凸包边界的有符号距离经平滑函数合成为像素 indicator。平滑度参数控制边/角的曲率，sharpness 参数控制从内部到外部的过渡宽度，因此同一表示可以从柔和扩散体连续调到近似硬的凸多边形。关键效率选择是直接在图像平面对投影点构造二维凸包，而不在每个视图先显式重建完整三维凸包。

CUDA rasterizer 对 tile 中的 smooth convex 计算该 indicator、颜色与 alpha，再按深度合成；点集、颜色、opacity、平滑/锐度参数均可由渲染误差反传。优化还根据梯度和投影覆盖控制 primitive 的增殖，使高误差区域可获得更多凸体。相较标准 3DGS，改的是几何 support 而非只改颜色核；代价是凸包投影、点集与 smooth indicator 的计算，并且单个凸原语无法直接表示非凸拓扑，仍需多个 primitive 组合。

% METHOD: 2412.04459
\paragraph{Sparse Voxels Rasterization（2024）\cite{2412_04459}}
Sparse Voxels Rasterization（SVR）验证实时显式 radiance rendering 不必以 Gaussian 为基本核。它以八叉树叶节点的非重叠稀疏体素表示场景，每个叶体素保存密度、尺度和外观（包括 SH）属性；对复杂区域继续细分、空区域不分配节点，因此几何 LOD 由空间层级而非 Gaussian 重叠堆叠提供。

渲染时体素投影到图像，而不是沿 ray 逐样本 march。由于叶体素在三维中不重叠，方法可用随 ray 方向变化的 Morton key 对投影体素进行排序，使同一 ray 的合成顺序与空间深度一致；随后以 tile 风格 rasterization 和 alpha 合成输出颜色。这种不交叠结构避免多个 Gaussian 排序互换产生的 popping，也使每个体素拥有明确的占据区域。

优化从粗层八叉树开始，根据图像重建误差逐步细分需要的叶节点，再优化其 density/SH 等参数；因而同时提供体素融合、pooling、marching cubes 等规则网格接口。训练仍将渲染图与多视图真值比较反传，但 primitive 的生长是八叉树细化而非 clone/split/prune。其边界在于高频薄结构可能需要大量细小体素，体素分辨率/内存和 tree traversal 会成为瓶颈；它解决的是显式体素的排序与渲染路径，并不等同于 Gaussian 的连续核抗锯齿特性。

% METHOD: 2501.03229
\paragraph{Gaussian Masked Autoencoders（2025）\cite{2501_03229}}
Gaussian Masked Autoencoders（GMAE）将 Gaussian 当作单张图像 MAE 的中层结构化 latent，而不是重建一个跨视图场景。普通 MAE 从未遮挡 patch 编码后直接预测被遮挡 RGB patch，语义特征虽强，却没有显式的层、尺度或深度组织。GMAE要求 decoder 先构造一组可 splat 的三维 Gaussian，再把它们渲回二维图像完成同样的 masked reconstruction。

图像被 patchify 并随机遮挡，ViT encoder 只处理可见 token；decoder 将可见 latent、mask token和可学习查询变换为每个 Gaussian 的中心、尺度、旋转、颜色、opacity等参数。高斯在单图相机坐标内可沿深度移动、大小可变，随后由可微 splatting 合成重建图，并只在 MAE 的像素重建目标上端到端训练。因此模型并未得到多视角真几何，却被迫以可遮挡、可分层、空间非均匀的 primitive 集解释图像。

这种中间表示的价值在于其非均匀分配：复杂物体边界与高信息区可获得更密/更小的 Gaussian，低信息区用更少/更大的 Gaussian；深度顺序可成为 figure-ground、层化和边缘的零样本线索。训练目标和下游仍是二维自监督表示学习，不能把输出中心当作标定相机下的可靠三维重建。与 3DGS 的共同点是 differentiable splatting，差异是没有 per-scene 多视图优化、也不以新视图合成为目标。

% METHOD: 2501.18630
\paragraph{Deformable Beta Splatting（2025）\cite{2501_18630}}
Deformable Beta Splatting（DBS）从核函数支持域切入。Gaussian 的指数尾部理论上无限，实际 rasterizer 必须人为截断；线性有限核虽避免这一点，形状却固定。DBS以可学习形状参数的 Beta 分布作为 splat 核，在有限 support 内让每个 primitive 在平缓覆盖与尖锐边界之间调节。

空间部分将投影坐标按基元尺度归一化后送入 Beta 型函数；形状参数改变中心集中度、边缘衰减和各向变形，使同样的有限椭圆支持可适配平面与细节。渲染仍按基元投影、核值乘 opacity、前到后合成进行，但核的零点在边界内被明确限定，tile 不再依赖高斯尾部阈值决定是否相交。论文还将 Beta 结构用于外观编码，区分更平滑的漫反射变化和更集中的镜面变化，而非只以低阶 SH 承担所有频率。

在密度控制上，方法讨论核分布改变后如何保持基元重定位/采样的分布一致性：通过相应的 opacity 正则与归一化关系，MCMC 式复制/重分配不依赖高斯的特定积分性质。优化仍由图像重建损失和正则驱动位置、尺度、Beta形状、颜色与opacity。有限核降低无效远处贡献却要求合适边界/形状学习；若形状参数退化或支持不足，细缝与半透明过渡可能出现缺口。

% METHOD: 2503.10148
\paragraph{3D Student Splatting and Scooping（2025）\cite{2503_10148}}
该方法把 3DGS 视为未归一化混合模型，并同时改变组件分布与混合符号。Student-$t$ 基元除位置、协方差、颜色和 opacity 外还有自由度 $\nu$：$\nu$ 小时是厚尾、接近 Cauchy 的广域组件，$\nu\to\infty$ 时回到 Gaussian。因此同一场可让大而平坦区域选厚尾、小而复杂区域选较尖分布，而不是所有基元共享指数衰减。

Splatting and Scooping进一步允许 opacity 通过 $\tanh$ 取正或负值。正组件向混合加入密度/颜色；负组件从已有混合中扣除贡献，可由一对共位组件表示环孔等“先覆盖、再挖掉”的非单调拓扑。论文不采用平方混合等会产生所有组件两两交互的形式，仍维持类似前到后 compositing 的线性复杂度，但以能量形式处理带负权的模型使学习目标有定义。

正负分布、自由度、协方差和opacity相互耦合后，普通 SGD 容易不稳定。方法将损失加上 opacity 与尺度正则定义后验能量，以 stochastic-gradient Hamiltonian Monte Carlo 采样参数；随机梯度噪声通过动量/摩擦项处理。低重要度组件会被重定位到高重要度组件附近并按 Student-$t$ 关系调整参数，新增零 opacity 组件再回收使用，替代 vanilla clone/split。重建监督仍是 $L_1$--D-SSIM。其额外表达也带来负 opacity 的稳定性、SGHMC超参数和对称 Student-$t$ 形状的限制，不能任意表示非对称核。

% METHOD: 2503.19232
\paragraph{HoGS（2025）\cite{2503_19232}}
HoGS处理同一场景同时含近处细节和极远目标时 Cartesian 3DGS 的参数病态：把高斯中心往很远处移动需反复积累巨大坐标更新，而位置变化还会改变投影尺度，令“远处但屏幕尺寸稳定”的目标难以优化。它把位置和尺度一起重参数化到齐次射影坐标，而不是再为远景单独添加 sky sphere 或背景模型。

一个 Cartesian 中心 $\mu$ 写成带权四维向量 $\tilde\mu=[\tilde x,\tilde y,\tilde z,w]^T\propto[\mu,1]^T$，回到 Euclidean 空间时除以 $w$。关键是尺度也共享同一个齐次分量：$\tilde s=[\tilde s_1,\tilde s_2,\tilde s_3,w]^T$，再经同样除法得到实际尺度。因而在固定前三个齐次分量时，较小的 $w$ 同时对应更远的位置与相应缩放，在靠近原点的相机下保持近似相同的屏幕投影；$w=0$ 还能表述射影无穷远方向。

渲染前将齐次位置和尺度还原为普通三维中心/协方差，再走原始投影、SH、alpha 合成与 $L_1$--D-SSIM 优化。适应这一表示后，论文关闭 vanilla 3DGS 定期删除 world-space 大高斯的规则，避免远景 primitive 被误当成异常尺度裁掉；其他致密化和渲染流程沿用基线。它改善的是深度跨度的优化坐标条件，而不会凭空解决远处被遮挡或训练图未覆盖的几何，且同一原点/投影设置是其参数化先验的一部分。

% METHOD: 2603.18707
\paragraph{From ex(p) to poly（以有限支撑多项式替换指数核）（2026.03）}
From ex(p) to poly~\cite{2603_18707} 不改变已有 3DGS 场景的高斯中心、协方差、颜色或 alpha 合成，而是替换每个投影二次型 $x=Q_i(v)$ 上的核求值。原始 splat 使用 $g(x)=\exp(-x/2)$，指数函数计算昂贵，虽在工程中会截断却没有真正的紧支撑。论文以 ReLU 激活的 $N$ 阶多项式 $f_N(x)=\max(\sum_{j=0}^N c_jx^j,0)$ 近似该核，使超出零点的贡献严格为零。

\textbf{如何得到并使用核。} 系数不是在任意一维采样上拟合：作者将二维屏幕空间均匀面积采样映射到二次型区间，再在原 3DGS 截断阈值所覆盖的区间对指数函数做 $\ell_1$ 拟合。运行时光栅化器以多项式值代替指数值，原有投影、深度排序和 alpha 混合保持不变。一阶形式计算最少；高阶可更贴近指数，但二次项可能在目标区间外重新为正，需在首个根处截断，而单调三次项则避免这一问题。

\textbf{紧剔除和抗混叠。} 线性核的零点给出与不透明度无关、可预计算的包围半径；使用阈值时还可得到依赖 alpha 的更紧半径。tile culling 因而只枚举真正可能有贡献的区域。论文也重新推导核积分所需的归一化，使协方差低通滤波在多项式核下仍保持跨视图的正确强度。它是对既有 3DGS 资产和渲染管线的可替换计算部件，而非重新优化一套场景表示；近似阶数、截断和滤波归一化需共同选择，不能仅把指数调用机械替换掉。

% METHOD: 2604.24994
\paragraph{Power Foam（有界 power diagram 的统一光追与光栅化）（2026.04）}
Power Foam~\cite{2604_24994} 并非把高斯替换为另一种重叠核，而是以不重叠的凸单元分割空间。Radiant Foam 的 Voronoi 单元可沿相邻面以近似常数代价行走，适合光线追踪，但单元可能无界，难以给 tile-based rasterizer 一个紧的屏幕包围。论文把 Voronoi 图推广为带权的有界 power diagram，使同一表示能在需要复杂光传输时追踪、在实时显示时光栅化。

\textbf{带权单元和可微边界。} 每个站点 $p_i$ 除位置外学习平方半径/权重 $r_i^2$，单元由最小化 $\|x-p_i\|^2-r_i^2$ 的区域定义，并与相应球界相交而受限。权重同时影响相邻单元的面和球边界，故即使界限与单元尺度不匹配也有梯度；不相交单元不会产生远处的伪相邻面。单元可携带密度与方向辐射以做体渲染；面向表面的版本显式区分内外界面，并把可微纹理附着到这些界面以解耦几何和外观。

\textbf{两条渲染路径。} 光追沿 power diagram 的对偶邻接关系逐单元前进，并把球界纳入相交长度计算。光栅化则按站点深度全局排序，投影每个有界单元的可见贡献并合成；论文利用分区的结构性质避免普通半透明 splat 深度启发式造成的 popping。表示的目标是统一两种枚举方式，而非声称它们有相同成本：power cell 的构建、权重优化和单元复杂度仍是额外约束，且它与独立高斯的重叠表示在数据结构上不同。

% METHOD: 2407.09733
\paragraph{Textured-GS（2024）\cite{2407_09733}}
Textured-GS改变单个 Gaussian 的外观坐标。原 3DGS 用相机到高斯中心的视线方向评估 SH，因此同一视角下该高斯覆盖的所有像素得到同一种颜色；一个较大的 Gaussian 无法表示其椭球表面上的棋盘、边界或 opacity 变化，只能靠增加更多基元。

方法把 Gaussian 看成带尺度和旋转的椭球。对每条像素 ray，先在椭球局部坐标中解二次方程得到最近的 ray--ellipsoid 交点，再按各轴尺度归一化为表面参数 $n$。颜色不再是 $Y_{lm}(v)$ 的函数，而用同样数量的 SH 系数评估 $Y_{lm}(n)$；因为 $n$ 随同视图中不同像素而变，SH系数成为椭球表面的空间纹理。它还给每个 SH 系数增加 opacity 通道，使局部 opacity 可变并改变感知形状，而不仅纹理化颜色。

投影与前到后 alpha 合成仍按 3DGS 进行，但 rasterizer 需额外计算 ray--ellipsoid 交点并在每像素按局部参数评估颜色/opacity。论文以 Mini-Splatting 的较紧凑高斯集初始化，冻结其位置、尺度和旋转并移除 adaptive densification，专门学习纹理 SH；渲染损失仍是 $L_1$ 与 D-SSIM。由于 SH 输入每像素不同，纹理与 opacity 系数的梯度也需按像素增量累积，而非将一个视角统一汇总。该方法以额外交点计算和 SH/opacity存储换取单基元更强的局部表达；它对椭球表面参数化的近似敏感，并会增加训练与渲染开销。

% METHOD: 2411.18625
\paragraph{Textured Gaussians（2024）\cite{2411_18625}}
Textured Gaussians为每个 Gaussian 加一张固定分辨率二维纹理图，而不是只以 SH 在整个 primitive 上给出统一颜色。原始高斯的空间变化只来自椭圆核的衰减，单标量 opacity 也只能形成椭球轮廓；方法希望一个基元同时携带高频局部纹理与非椭球的透明边界。

纹理贴在 Gaussian 两条最大主轴张成的局部平面上。渲染某个像素时，从相机原点向该像素发射 ray，求它与此平面的交点；交点沿两条主轴的坐标按各自 scale 和纹理覆盖系数映射到 $(u,v)$，并用双线性插值读 RGB 与 alpha。最终颜色为原有空间常量、可随视角变化的 SH 基色加上纹理 RGB；最终 alpha 为纹理 alpha、Gaussian 核值与基元 opacity 的乘积。于是 alpha-only map 可以改形状，RGB/RGBA map 可以在同一 Gaussian 内表达空间细节。

前到后合成形式保持不变，只是每个贡献项现在是 $c_i(u,v)\alpha_i(u,v)$。论文先训练普通 3DGS，再以其参数和逐 Gaussian texture map 联合优化第二阶段，并停止 ADC 以保持基元数可比较；监督仍是 $L_1$ 与 SSIM。自定义 CUDA 实现 ray--plane intersection、UV 与纹理采样。它比 SH 表面函数更接近显式 texture atlas，但每 Gaussian 的纹理显存、分辨率和 ray--plane 近似都会影响速度和质量，且大视差下平面纹理并不能完全描述三维体积外观。

% METHOD: 2507.20200
\paragraph{Neural Shell Texture Splatting（2025）\cite{2507_20200}}
Neural Shell Texture Splatting将几何与外观彻底拆开：Gaussian（论文以 2D oriented splat 为几何底座）只存中心、旋转、尺度和 opacity，移除逐基元 SH；高频外观改由一个全局、多分辨率 hash-encoded shell texture 提供。动机是纹理复杂区域若仍由每个高斯携带 SH，优化会靠不断 split 产生大量针状 primitive，把几何密度和纹理频率混为一个问题。

对每个 pixel--splat 精确交点 $(u,v)$，以 splat-to-world 齐次变换直接得到世界坐标 $p_u$，查询多层 hash grid 的特征。不同深度的这些特征仍按 alpha/transmittance 先合成为屏幕 feature map，而不是逐 Gaussian 立即解码 RGB；一个小 MLP 在 deferred rendering 阶段将合成特征解码为最终颜色。这使同一连续纹理场可跨多个 primitive 共享，而高斯只决定表面覆盖与可见性。

训练从 2DGS 初始化，按 coarse-to-fine 激活 hash 编码层以先拟合低频再引入细节；损失包括标准 RGB $L_1$--SSIM、2DGS 的深度与法线正则以及 alpha map 项。交点到 world 的直接变换让纹理梯度同时回传至 shell features 和 splat 几何。输出可从全局纹理场取样到网格，避免逐高斯 atlas 的拼接。它以额外 hash volume/MLP 代替 SH 内存和过度致密化，但受 shell 交点准确性、哈希碰撞和全局纹理对未观测区域的泛化限制。

% METHOD: 2511.14357
\paragraph{IBGS（2025）\cite{2511_14357}}
IBGS将 3DGS 渲染颜色拆成自包含的 SH base color 与从训练图像取回的 image-based residual。SH 可高效概括低频及一部分视角相关外观，但复杂纹理、镜面反射和细粒度颜色若全压进逐高斯系数会增加模型容量。IBGS选择在新视图渲染时再次利用邻近源图像，而不是给每个 Gaussian 分配大型纹理图。

对目标像素，标准 3DGS先渲出 base color、深度/几何信息；根据该点的可见表面位置，将它重投影到多个选择的训练视图，采样对应源像素并以可见性/几何关系形成 warped colors。论文还提取源图像特征并处理不同视图的曝光，使各来源可比较。一个 residual prediction module 将 base color、warped color、视角及局部特征融合，输出 $\Delta c$，最终颜色为 $c_{base}+\Delta c$。因此残差可从真正拍摄过的相邻视图借回高频证据，而 base 3DGS提供新视图的几何锚点。

训练时从多视图重建误差反传，既优化 Gaussian 场，也优化特征/残差与曝光校正模块；渲染接口需要访问源图和相机，而非只加载单个 PLY 模型。它不同于全局 neural texture 或逐 Gaussian atlas，优势是保留原始照片细节，边界则是源视图覆盖、重投影深度和曝光一致性：遮挡、误几何或目标视图远离全部训练相机时，warp color 可能错误，系统需回退到基础颜色而不能保证细节真实。

% METHOD: 2511.18873
\paragraph{Neural Texture Splatting（共享神经局部纹理）（2025.11）}
Neural Texture Splatting~\cite{2511_18873} 要解决的是普通 3DGS 中单个高斯只提供空间常量颜色与不透明度，因而难以在较少基元下表达高频纹理的问题。直接给每个高斯存一张 RGBA 纹理虽能提高局部表达力，却会随高斯数线性增加存储；为每个高斯配一个网络也会使查询昂贵。该方法保留高斯的位置、旋转、尺度及前到后的 alpha 合成，但把每个高斯的外观补偿改为一个由共享网络生成的局部纹理场。

\textbf{局部纹理查询与渲染。} 对一条相机光线与第 $k$ 个高斯相交得到的世界点 $\hat{x}$，先用该高斯的中心 $\mu_k$、旋转 $R_k$ 和尺度 $S_k$ 变换到局部坐标 $x_k=S_k^{-1}R_k(\hat{x}-\mu_k)$。局部纹理场在 $x_k$ 查询 RGBA 值：RGB 作为基色的附加颜色，alpha 作为该高斯不透明度的附加项，再进入原有的深度排序与 alpha 合成。这样纹理是在高斯自身坐标系中附着的；查询点越出局部包围盒便不贡献，避免一个基元的纹理泄漏到无关空间。

\textbf{从显式纹理到共享神经纹理。} 论文先说明显式版本：每个高斯携带三个局部 RGBA 平面 $F_{xy}^k,F_{xz}^k,F_{yz}^k$，在三个投影面上双线性采样并取平均。完整 NTS 不逐个存储这些平面，而是维护全局三平面特征；在高斯中心读取特征 $g_k$，把 $g_k$、中心位置和视线方向（动态数据还加入时间）送入解码器，预测该高斯的局部三平面。颜色和不透明度使用分开的全局平面与解码器，因而视线方向只进入颜色分支。为控制局部平面的预测成本，论文再以两个一维向量的外积来因子化平面 RGBA 值，而非直接输出完整二维格。

\textbf{优化范围与取舍。} 训练仍以渲染图像和观测图像的重建项为主体，并对三平面 RGBA 值施加 $\ell_1$ 稀疏约束；射线--高斯相交为纹理查询提供位置。它与逐高斯显式纹理的差别不在于把纹理简单搬到全局网格，而在于“共享全局先验预测局部、随高斯变换而移动的纹理”。这种共享降低了重复存储，也让相邻高斯可借用共同特征；代价是全局平面的分辨率及解码器容量会限制可恢复细节，并增加每次渲染的查询计算。

% METHOD: 2604.01204
\paragraph{Neural Harmonic Textures（原语附着的谐波特征与延迟解码）（2026.04）}
Neural Harmonic Textures~\cite{2604_01204} 针对单个 3DGS/2DGS 等拉格朗日原语内部颜色近似恒定的问题。单纯增加原语数可补高频纹理，却同时增加存储与光栅化负担；若在每个射线--原语交点运行 MLP，又使计算随交点数增长。论文将局部特征锚定在随原语运动的虚拟支架上，再把这些特征转为可在图像空间一次性解码的谐波分量。

\textbf{支架与谐波纹理。} 对每个原语，在其周围定义虚拟四面体并在四个顶点存潜特征。光线在原语最大响应点的局部坐标给出重心权重，先插值得到特征 $f_i$；随后对其施加 $\sin(f_i)$ 与 $\cos(f_i)$。顶点特征差异会调制局部频率，而该原语的核加权不透明度成为谐波振幅。支架随原语坐标一起变换，所以它保持显式原语在形变、编辑或不同 primitive 类型上的局部性，而不需要全局规则网格。

\textbf{从交点解码转为 deferred decoding。} 沿一条光线，系统仍按 transmittance 和 alpha 将各原语的谐波向量累积；但不在每个交点解码颜色。累积向量与二阶 SH 编码的视线方向在图像空间输入一个浅层 MLP，一次输出像素颜色。因而高维空间变化信号在合成前被保留，神经解码的次数按像素而非射线--原语交点计。该模块可以接到 3DGS、2DGS 或 triangle splatting 的可微合成之后；代价是每原语增加支架特征，且最终外观仍依赖延迟网络的容量和训练分布。

% METHOD: 2410.04974
\paragraph{6DGS（2024）\cite{2410_04974}}
6DGS把观察方向纳入 primitive 的概率表示，而不是只将方向送入颜色 SH。标准 3DGS有三维位置协方差和固定 opacity，SH 只能修改外观；在烟、云、半透明或强视角依赖材质中，同一空间高斯的可见性本身也应随方向改变。方法为每个基元建立空间位置与三维方向的联合六维 Gaussian。

给定相机方向 $d$，联合分布通过条件 Gaussian 公式切片为一个条件的三维空间 Gaussian：空间均值由空间--方向协方差回归项偏移，条件协方差由 Schur complement 得到。方向分量相对于均值方向的 Mahalanobis 距离再产生条件因子 $f_{cond}=\exp(-\lambda_{opa}D)$，与基础 opacity 相乘；$\lambda_{opa}$ 可固定或逐 Gaussian 学习，控制方向对可见性的影响。颜色则保留 SH，而不是退回为每方向独立 RGB。

渲染时针对当前视线把每个 6D 基元切为相应的条件 3D Gaussian，再走标准投影、tile 排序与 alpha 合成。密度控制同样可访问方向条件信息：基元是否需要增殖或裁剪不只依赖空间误差，也反映其对不同观察方向的有效支持。训练仍由多视图渲染损失优化六维均值/协方差、opacity和SH。方向维增加了对体积与视角相关现象的表达，却也带来条件矩阵运算、更多参数与切片开销；方向没有强变化的场景未必需要它。

% METHOD: 2512.14180
\paragraph{Spherical Voronoi（以可微球面分区表示方向外观）（2025.12）}
Spherical Voronoi~\cite{2512_14180} 替换的是 3DGS 高斯颜色中的球谐方向函数，而不是空间高斯的中心、协方差或 alpha 合成。有限阶 SH 对平滑方向变化有效，却难以紧凑地表达局部尖锐的方向信号；多个球面高斯核又会重叠竞争、使优化依赖初始化。论文把单位球面的方向域分为可学习站点控制的软 Voronoi 区域，以连续权重输出颜色或其他方向量。

\textbf{球面函数。} 对 $K$ 个单位方向站点 $s_k$、关联值 $c_k$ 与正温度 $\tau_k$，方向 $\omega$ 的输出为 $f_{\mathrm{SV}}(\omega)=\sum_k w_kc_k$，其中 $w_k=\operatorname{softmax}_k(\tau_k s_k^\top\omega)$。内积衡量球面邻近性；$\tau_k$ 小时边界平滑，增大时趋近硬 Voronoi 分区。每个高斯保存一组 $(s_k,c_k,\tau_k)$ 来替代其 SH 系数，渲染时依据视线方向求该函数得到颜色，余下的投影、排序和 alpha 合成仍沿用 3DGS。因此它把“外观如何随方向变化”与“高斯在空间中如何分布”明确解耦。

\textbf{反射版本。} 论文还在 2DGS 的 deferred rendering 管线上应用这一方向表示。几何 pass 输出位置、法线、反射率和深度；每个原语额外学习漫反射、粗糙度及 SV 参数。镜面项分为远场与近场：远场由反射方向查询立方图，近场则由邻近空间光探针的 SV 函数按距离混合；粗糙度调节分区温度，使低粗糙度可保留更尖锐的方向峰。最终将该镜面项与漫反射合成，而非把反射硬编码进单一 SH 系数。

\textbf{适用边界。} 这种参数化提供了局部可控的方向分区，并通过 softmax 保持站点的可微竞争，但站点数、每站点值及温度仍会增加每个高斯的外观参数。它解决的是方向外观的表达与优化问题，不自动补足全局光照、未观测反射或空间几何的证据。

% METHOD: 2501.15619
\paragraph{GaussianToken（2025）\cite{2501_15619}}
GaussianToken把二维 Gaussian 变成图像 tokenizer 的基本 token。离散 VQ tokenizer 通常把固定网格位置的 feature 向量替换为 codebook 条目，token 的几何 support 未被编码；GaussianToken将每个 token 表示为二维位置、旋转、两轴尺度和特征向量，让“是什么”与“覆盖在哪里、覆盖多大”分开处理。

编码器先输出图像特征，再预测一组 Gaussian 几何参数及特征。只有特征向量进入 VQ codebook 做最近邻离散化，中心、尺度与旋转保持连续；解码时把量化 feature 与连续几何一起用二维 Gaussian splatting 写回特征图，随后由图像 decoder 生成 RGB。几何参数决定每个离散码字向邻域传播的 footprint，故多个 token 可在连续平面上移动、重叠或改变形状，而不受固定 patch 栅格限制。

训练包含图像重建、codebook/commitment 等 tokenizer 约束并通过 splatting 反传；使用时输出的是适合生成模型的离散 token 加连续布局，而非三维高斯场。相对直接把 2D Gaussian 当图像压缩器，它的核心接口是“可量化语义特征 + 不量化的连续几何”；表达力和压缩率取决于 codebook 大小、token 数与几何参数开销，连续布局也不能自动提供跨图真实几何对应。

% METHOD: 2512.19108
\paragraph{GaussianImage++（失真驱动的二维高斯图像表示）（2025.12）}
GaussianImage++~\cite{2512_19108} 将 Gaussian splatting 用于单幅二维图像表示：每个基元只有图像平面位置、二维协方差和颜色，不需要相机投影、深度排序或球谐颜色。其核心问题是把 3DGS 的位置梯度式 densification 直接移到二维图像时，梯度累积不足以可靠决定在哪里添高斯。论文因此以当前图像重建失真而非位置梯度来分配有限的基元预算。

\textbf{D$^3$：由像素误差决定生长。} 表示从稀疏二维高斯开始，在固定优化间隔计算原图与当前渲染之间的逐像素失真。调度器在剩余预算中选择一批高误差像素，在这些像素位置添加新高斯，并以该像素的目标颜色初始化其颜色；新增数量随剩余预算递减，以免一次耗尽总数。协方差直接作为二维矩阵优化。直接优化可能产生非正半定矩阵，因而优化过程中检查并剪除无效基元。这样，平坦区域不会与持续高误差的边缘或纹理竞争同样数量的高斯。

\textbf{CAF：让早期覆盖和后期细节使用不同滤波尺度。} 单一固定低通滤波无法同时处理稀疏早期的大面积空洞与后期的细小结构。Content-Aware Filtering 为新加入的高斯赋予随训练阶段变化的滤波强度：早期加入者的有效协方差经较强滤波而覆盖更大区域，后期新增者以较弱滤波专注残余细节；已有基元保留其加入时的滤波状态。该规则与 D$^3$ 配合，前者控制新增基元的位置，后者控制其初期覆盖范围。

\textbf{压缩接口与边界。} 论文还把位置、协方差与颜色分开处理，采用可学习的标量量化和量化感知训练，使拟合出的二维高斯集合可作为压缩表示的一部分。这里的“压缩”建立在单图二维信号上，不能等同于三维场景的几何压缩；D$^3$ 也仍是逐图优化流程，其结果与预算、损失图和优化设定相关。

% METHOD: 2312.11461
\paragraph{GAvatar（2023）\cite{2312_11461}}
GAvatar从文本生成可动画的 Gaussian avatar。直接在世界坐标独立优化大量高斯，会使姿态变化时缺少一致变形约束，而 SDS 这类高方差生成监督又容易让邻近高斯学到不协调的颜色、尺度和旋转。方法以可随人体姿态形变的 primitive 为骨架，将每个 Gaussian 绑定到一个局部 primitive 坐标系。

每个表面 primitive $V_k=\{P_k,R_k,S_k\}$ 关联一组局部高斯 $(p_k^i,r_k^i,s_k^i,c_k^i,\sigma_k^i)$。给定姿态，primitive 的平移、旋转和尺度把局部中心、尺度与旋转变换到世界空间，因此同一组高斯自然随人体运动。为了不逐个直接学习属性，方法在 canonical rest pose 的位置查询两个 neural implicit field：$\mathcal H_\phi$ 输出颜色、旋转与尺度，$\mathcal O_\psi$ 输出 opacity。共享连续场使局部邻域有一致属性，也将属性跨姿态 canonicalize；位置正则限制高斯不要漂离所属 primitive。

几何并不只依赖透明 Gaussian 云。GAvatar以 SDF 表示表面，并把 opacity 参数化为 SDF 值经 NeuS 风格钟形核后的结果，使高 opacity 集中在零水平集附近。SDF允许 DMTet 提取网格，并加入 Eikonal、网格 alpha 与 normal consistency 正则；提取网格的 normal/RGB 还可接受 SDS 监督和纹理微调。总训练目标结合 RGB SDS、位置正则、SDF/网格正则和 normal SDS，优化局部高斯、隐式属性/SDF、primitive 运动校正和人体形状参数。最终 RGB 仍由 Gaussian splatting 生成，但输出同时可动画并可导出带纹理网格；效果受文本生成先验与人体 primitive 模型限制，不能保证任意服装或拓扑的物理真实性。

% METHOD: 2403.14627
\paragraph{MVSplat（2024）\cite{2403_14627}}
MVSplat是从少量已知相机图像直接前馈预测 Gaussian 场的模型，避开每场景数万步优化。它将每个输入像素对应一个高斯，关键在于先得到可靠深度：若中心位置错误，即使 opacity、协方差和 SH 颜色预测正确也无法合成稳定新视图。

输入 $K$ 张图像和投影矩阵后，浅层 CNN 与含 self/cross attention 的 Transformer 提取跨视图特征。对每个参考视图，在 inverse-depth 上采样 $D$ 个候选深度，把其他视图特征按相机几何 plane-sweep 回投影，计算点积相关并对其他视图平均，得到 $H/4\times W/4\times D$ cost volume。轻量二维 U-Net 用参考特征预测 cost-volume residual，并在低分辨率层交换各视图的注意力；这保持网络可接受任意数量的输入视图而不使用三维卷积。归一化后的深度维分数与候选深度加权平均得到深度图，另一个二维 refinement 网络进一步修正它。

每个预测深度按相机反投影为 Gaussian 中心；并行的参数分支预测 opacity、协方差/尺度旋转和 SH 颜色，形成 pixel-aligned 的 $H\times W\times K$ 高斯集合。网络对新视图用普通可微 alpha splatting 渲染，只以图像重建损失端到端训练，不依赖显式深度真值。因而多视图匹配提供几何、参数分支提供外观，随后即可一次前向输出场并实时渲染。它的限制正来自这些前提：误匹配、遮挡或极稀疏视图会污染 depth/cost volume，且固定的每像素基元预算可能在无纹理或极细结构处不足。

% METHOD: 2509.20207
\paragraph{PU-Gaussian（2025）\cite{2509_20207}}
PU-Gaussian将各输入点附近的表面视为一个可采样的各向异性 Gaussian，用于点云上采样，不是用点云训练新视图 renderer。给定稀疏点集 $\mathcal X$，目标是生成 $rN$ 个均匀点，同时保留边缘和局部曲面；传统上采样网络直接回归每个新点，缺少对局部不确定性和方向性的显式模型。

Gaussian prediction network 先将输入坐标投到特征空间，再以 Point Transformer 取得局部/全局描述；三条 head 分别预测尺度、四元数旋转和相对输入点的均值 offset，形成每点 $\mathcal N(\mu_i,\Sigma_i)$。从每个分布用 reparameterization trick 抽取 $r$ 个粗点，并丢弃超过两倍标准差的样本。随后将重复的原点特征与粗点坐标拼接，再经第二个 Point Transformer 和 offset head 预测残差，得到边缘更清晰的 refined 点云。

训练第一阶段以 Chamfer distance 和 Gaussian Mahalanobis 正则先学习局部分布，使真值点接近其对应预测均值；第二阶段端到端优化 refined 输出的 Chamfer 距离。推理中 refinement 可再执行一次，利用抽样随机性获得较稳定的密集点集。Gaussian 在此只定义三维采样密度，没有颜色、opacity、投影或 alpha compositing；其限制来自各向 Gaussian 的局部平滑假设，极尖锐角、稀疏缺口及随机采样方差仍需靠 refinement 纠正。

% METHOD: 2512.10685
\paragraph{Sharp Monocular View Synthesis in Less Than a Second（单图前馈高斯重建）（2025.12）}
Sharp Monocular View Synthesis in Less Than a Second~\cite{2512_10685} 针对单张 RGB 图像直接生成可新视角渲染的高斯场。常规 3DGS 依赖多视图和逐场景优化；单图任务则同时面对不可见表面、遮挡与尺度歧义。论文把这一过程改成一次前馈预测：网络从输入图像产生两层深度及每个像素对应的高斯属性，再按普通高斯光栅化器从目标相机渲染。

\textbf{两层深度初始化。} 输入是 $3\times H\times W$ 的 RGB 图像。以 Depth Pro 为编码器，作者复制末层卷积以输出两通道深度，并微调解码器及低分辨率编码器层，使其服务于新视角任务。训练时由预测深度和真值深度构造缩放图，对深度进行调整；两层而非单层深度的目的，是让同一像素容纳前景背后的潜在表面，从而对遮挡或轮廓处的单图歧义留出表示空间。

\textbf{从深度图到高斯。} 初始化器下采样图像与两层深度，将像素与相应深度反投影为中心位置；颜色取自输入图像，旋转先设为单位旋转，不透明度先设为常数，尺度随深度设定。随后 Gaussian decoder 预测对位置、尺度、旋转、颜色和不透明度的修正，最终得到每层每像素一个高斯。论文的输出张量为每个高斯 14 个属性通道，位置、尺度、四元数旋转、RGB 与不透明度均被显式预测；它不以球谐函数承担外观，而以输入视图颜色及其预测修正初始化。

\textbf{训练约束与渲染。} 系统在输入视图与新视图上比较渲染 RGB，并在新视图加入感知项；alpha 的二元交叉熵、首层深度的 $\ell_1$ 项、第二层深度的总变分项共同约束两层结构。为压制高视差梯度附近的漂浮高斯，论文还给出 floater 正则，并限制位置的图像平面偏移及投影方差范围。推理时只需执行深度预测、属性解码和标准 splatting，不再进行逐图迭代优化。它的速度来自前馈化而非消除几何不确定性：单张图中不可见区域仍依赖训练数据先验，两层深度也不能保证恢复真实的完整三维拓扑。

% METHOD: 2606.24876
\paragraph{FLAT（前馈 latent 三角形 splatting）（2026.06）}
FLAT~\cite{2606_24876} 面向单图生成可探索场景时，选择表面对齐的 triangle splat 而非体积 3D Gaussian。三角形有明确表面，便于后续图形与仿真接口；但其朝向一旦预测错误就几乎没有屏幕贡献，直接回归三维顶点或全局旋转会造成退化和梯度消失。论文把冻结的视频扩散先验与前馈场景解码器结合，并专门设计三角形参数化和覆盖函数来稳定这个回归任务。

\textbf{条件 latent 到三角形。} 输入单张 RGB 与目标相机轨迹。点云控制视频和相机嵌入条件化冻结的视频扩散模型，得到去噪视频 latent；复用预训练视频 VAE 的 RGB decoder 骨干，在中间特征上接输出头预测三角形属性。相机以逐像素射线的参考点--方向表示编码，经零初始化注入块与视觉 latent 对齐。每个 token 对一个局部图像区域预测深度、二维形状、旋转残差、颜色、不透明度和锐度，并在单次前向中转成世界坐标三角形集合。

\textbf{射线中心参数化与可微覆盖。} 三角形中心位于 anchor ray 的预测深度处；其局部二维形状从正三角形经下三角矩阵变换获得，对角项强制为正以保证面积，非对角项控制剪切。随后以射线切平面为基，通过两个倾斜残差与平面内旋转抬升到三维，避免直接预测自由四元数。渲染时投影三个顶点，对每条边的归一化有符号距离做带偏移的截断响应，并把三个响应相乘再以锐度幂控制边缘；该窗口在边界外仍保留非零导数，且让梯度传至三条边，而非 max-edge 的单一路径。三角形按深度前到后 alpha 合成。论文另给出轻量测试时细化以把 triangle soup 转为不透明资产；生成质量仍依赖视频先验对未见区域的补全，且三角形数量、方向误差和细化步骤是这一表示的实际代价。
